\section{Discussion}\label{sec:suggestion}

\subsection{Grace Hopper}
NVIDIA's Grace Hopper architecture introduces tight CPU–GPU integration via a coherent interconnect and hardware‐accelerated contexts. Our evaluation on a Grace‐Hopper node shows that \sys{}'s instrumentation overhead drops below 1.2 µs per probe and data‐exchange bandwidth increases by 40 \% compared to PCIe. Looking ahead, Hopper's preemptible kernels and multi‐instance GPU (MIG) support could further reduce probe latencies and enable safe, per‐tenant instrumentation domains. Exposing finer‐grained context‐switch controls in hardware would allow \sys{} to migrate persistent kernels across SM partitions dynamically, improving load balancing and utilization in streaming or low‐latency scenarios.

\subsection{CXL Memory Pool}
The advent of CXL.mem promises virtually unbounded memory pools shared across hosts and accelerators. In our CXL experiments, dynamic latency insertion highlighted the impact of link saturation and access‐pattern sensitivity on end‐to‐end performance. Hardware support for explicit memory‐domain hints (e.g., "hot" vs. "cold" data), flow‐control credits, or quality‐of‐service registers would allow runtime systems like \sys{} to place instrumentation buffers in the most appropriate domain automatically. Moreover, integrating CXL pools with GPU‐resident shared‐memory allocators could eliminate round‐trip latency spikes, enabling truly transparent offload of large eBPF map structures without oversubscribing local HBM capacity.

\subsection{Suggestions for Hardware Vendors}
To unlock the full potential of dynamic, on‐device instrumentation, we recommend the following hardware primitives:
\begin{itemize}[leftmargin=*,itemsep=0pt]
  \item \textbf{Hardware Shared Queue (Semaphore) for Shared Memory.}  A lightweight, on‐chip queue or semaphore primitive—distinct from full cache coherence—would let CPU and GPU efficiently synchronize access to shared data structures (e.g., eBPF maps) without expensive cache‐line invalidations. Such a queue could provide producer/consumer semantics in hardware, reducing software‐polling overhead and improving latency under high contention.
  \item \textbf{Native API for Safe Self‐Modifying Code.}  Expose a protected code‐patch API in the GPU driver or runtime that validates and installs new kernel code pages at runtime. With support for atomic patch points and safe rollback, JIT engines could inject or update PTX snippets without relying on speculative buffer emulation. This would simplify dynamic instrumentation frameworks and reduce the complexity of user‐space sandboxing layers.
\end{itemize}
By incorporating these features, future GPU and interconnect designs can better support adaptive observability frameworks like \sys{}, delivering low‐overhead, secure, and flexible instrumentation for emerging AI and HPC workloads.
